\section{The premise graph}\label{sec:graph}

A single document can be checked in one pass. A project consists of many documents, written over months, and the questions a researcher
asks of it are graph questions: \emph{What does this result rest on? If this assumption changes, what else is affected? Has anything
ever been computed from this premise?} QAM answers them with a local SQLite database.

\subsection{Indexing}

Indexing reads the exported records under a project's record directories into the database; the deliberately broken examples are excluded by name, and the record of this article lies outside the default directories. It recomputes and verifies every entry hash and record hash on the
way, and checks the hash of the source file where it is available. A record that fails this verification is skipped as a whole, with a message. For every entry
the database stores the edges to its premises, indexed in both directions. In the repository from which the examples are taken, indexing reports 81
hash-verified entries, among them the two entries of the example rule.

Re-indexing rebuilds the entries and edges in one transaction but leaves stored receipts, the calculation cache and recorded replacements
untouched. A receipt therefore remains valid evidence of what was computed at the time, even after the record has moved on.

\subsection{Questions as recursive queries}

Three queries cover the questions above. Each is a recursive query over the edge table.
\begin{itemize}[leftmargin=1.5em]
\item \emph{Premises.} Starting from the hash of an entry, follow the premise edges to their end and return every entry reached,
with its statement, status and frame. For the example rule the answer is the rule itself and the preparation, both
\st{conditional}. These are exactly its open conditions in the sense of Definition~\ref{def:open}.
\item \emph{Impact.} Starting from the hash of an entry, follow the edges in the other direction and return every entry that depends
on it, together with every stored receipt whose premise trace contains it, either as a resolved or as a missing premise. For the
preparation of the example, the answer is the rule and the stored receipts of the runs of Section~\ref{sec:requests} that reached it: the successful run and the two rejected ones.
\item \emph{Chains.} A catalogue of rules records, for each recipe, the declarations it requires and the quantity it produces. The
query matches outputs of one rule to requirements of another and suggests paths of two to five rules. These are suggestions: execution
still checks values, frames and premises. The interval rule currently has no registered downstream rule.
\end{itemize}

\subsection{Corrections without overwriting}

The record is append-only, so an error is corrected by a new entry, and the database links the two by recording a \emph{replacement}
from the old hash to the new one, with a stated reason (the two hashes are stored as given; that they belong to indexed entries is not checked). A trigger rejects replacements that would form a cycle. Old entries and old
receipts are not changed. The impact query then shows what was built on the superseded entry, so the correction can be propagated
deliberately rather than by search and replace.

\subsection{Checking readiness before running}

The rule catalogue can also be asked whether a calculation is ready to run. For the interval recipe, the catalogue lists two required
declarations (the preparation and the boost) and one optional declaration (the measurement). Given the two required ones, the database
answers that the recipe is ready for validation and that its rule is indexed; the concrete values and the full premise chain are then
checked by the executor at run time. The division of labour is deliberate: the database answers questions about structure, and only
the executor computes.
